%%
%% part2-book-components/ch06-chapters.tex
%%
%% Chapter 6: Chapters and Sections — How to organize a book
%% into parts, chapters, sections, subsections, and
%% subsubsections, and how to cross-reference them.
%%

\chapter{فصل و بخش}
\label{chap:chapters}

\section{چرا این موضوع مهم است؟}
\label{sec:chapters-why}

یک کتاب خوب، فقط مجموعه‌ای از متن نیست؛ بلکه
ساختاری منظم دارد که خواننده را در مسیر یادگیری
هدایت می‌کند. فصل‌ها، بخش‌ها و زیربخش‌ها،
ستون فقرات این ساختار هستند.

\lr{MatinBook} از یک سیستم عنوان‌بندی پنج‌سطحی
استفاده می‌کند که بر اساس مقیاس \lr{Boyer/Stewart}
طراحی شده است. در این فصل، با این سیستم آشنا
می‌شوید و یاد می‌گیرید که چگونه ساختار کتاب
خود را سازمان‌دهی کنید.

\mnote{مقیاس \lr{Boyer/Stewart} استاندارد کتاب‌های درسی دانشگاهی است و در انتشارات معتبری مانند \lr{Wiley} و \lr{Pearson} استفاده می‌شود.}

\section{سطوح عنوان‌بندی}
\label{sec:chapters-levels}

\lr{MatinBook} از پنج سطح عنوان‌بندی پشتیبانی
می‌کند:

\begin{enumerate}
    \item \textbf{بخش} (\cmd{part}): بزرگ‌ترین
    واحد ساختاری.
    \item \textbf{فصل} (\cmd{chapter}): واحد
    اصلی کتاب.
    \item \textbf{بخش} (\cmd{section}): زیرمجموعه‌ی
    فصل.
    \item \textbf{زیربخش} (\cmd{subsection}):
    زیرمجموعه‌ی بخش.
    \item \textbf{زیرزیربخش} (\cmd{subsubsection}):
    عمیق‌ترین سطح.
\end{enumerate}

\mnote{توجه کنید که «بخش» در فارسی برای هر دو \lr{part} و \lr{section} استفاده می‌شود. در این کتاب، برای \lr{part} از «بخش اصلی» و برای \lr{section} از «بخش» استفاده می‌کنیم.}

\section{بخش اصلی (\lr{part})}
\label{sec:chapters-part}

بخش اصلی، بزرگ‌ترین واحد ساختاری در یک کتاب
است. یک کتاب می‌تواند چند بخش اصلی داشته باشد
که هر یک شامل چند فصل است.

برای ایجاد بخش اصلی:

\begin{latin}
\begin{minted}{latex}
\part{|\pc{مبانی}|}
\label{part:fundamentals}
\end{minted}
\end{latin}

\mnote{بخش اصلی در فهرست مطالب ظاهر می‌شود ولی شماره‌ی فصل را تغییر نمی‌دهد.}

\subsection{شماره‌گذاری بخش اصلی}
\label{sec:chapters-part-numbering}

بخش‌های اصلی با اعداد رومی شماره‌گذاری می‌شوند:
بخش اول، بخش دوم، بخش سوم و...

\subsection{نمایش بخش اصلی}
\label{sec:chapters-part-display}

هر بخش اصلی، یک صفحه‌ی جداگانه دارد که فقط
شامل عنوان بخش است. این صفحه، بدون شماره
است و به عنوان یک جداکننده بین بخش‌ها عمل
می‌کند.

\mnote{اگر کتاب شما کوتاه است (کمتر از ۸ فصل)، نیازی به بخش اصلی ندارید.}

\section{فصل (\lr{chapter})}
\label{sec:chapters-chapter}

فصل، واحد اصلی کتاب است. هر فصل، یک موضوع
مشخص را پوشش می‌دهد و معمولاً ۱۵ تا ۲۵ صفحه
است.

برای ایجاد فصل:

\begin{latin}
\begin{minted}{latex}
\chapter{|\pc{مقدمه}|}
\label{chap:introduction}
\end{minted}
\end{latin}

\subsection{شماره‌گذاری فصل}
\label{sec:chapters-chapter-numbering}

فصل‌ها با اعداد عربی شماره‌گذاری می‌شوند: ۱، ۲،
۳ و... این شماره‌گذاری در \cmd{mainmatter}
شروع می‌شود و در \cmd{backmatter} ادامه می‌یابد.

\subsection{سبک نمایش فصل}
\label{sec:chapters-chapter-style}

فصل‌ها در \lr{MatinBook} با یک سبک خاص نمایش
داده می‌شوند:

\begin{itemize}
    \item شماره‌ی فصل با فونت \pnum{۲۴pt} و
    رنگ کم‌رنگ.
    \item عنوان فصل با فونت \pnum{۲۰pt}.
    \item یک خط افقی رنگی زیر عنوان.
\end{itemize}

\mnote{سبک فصل در فایل \lr{mb-headings.sty} تعریف شده است. برای شخصی‌سازی، این فایل را ویرایش کنید.}

\subsection{فصل‌های بدون شماره}
\label{sec:chapters-chapter-star}

برای فصل‌های بدون شماره، از \cmd{chapter*}
استفاده می‌کنیم. این دستور برای پیشگفتار،
پیوست‌های بدون شماره و بخش‌های ویژه مناسب
است:

\begin{latin}
\begin{minted}{latex}
\chapter*{|\pc{پیشگفتار}|}
\addcontentsline{toc}{chapter}{|\pc{پیشگفتار}|}
\end{minted}
\end{latin}

\mnote{دستور \lr{\textbackslash addcontentsline} باعث می‌شود که این فصل در فهرست مطالب ظاهر شود.}

\section{بخش (\lr{section})}
\label{sec:chapters-section}

بخش، زیرمجموعه‌ی فصل است. هر فصل معمولاً
شامل ۳ تا ۵ بخش است. برای ایجاد بخش:

\begin{latin}
\begin{minted}{latex}
\section{|\pc{مقدمه}|}
\label{sec:introduction}
\end{minted}
\end{latin}

\subsection{سبک نمایش بخش}
\label{sec:chapters-section-style}

بخش‌ها در \lr{MatinBook} با یک سبک خاص نمایش
داده می‌شوند:

\begin{itemize}
    \item شماره‌ی بخش در یک \lr{colorbox}
    با پس‌زمینه‌ی آبی روشن.
    \item عنوان بخش با فونت \pnum{۱۵pt}.
\end{itemize}

\mnote{شماره‌ی بخش به صورت «فصل.بخش» نمایش داده می‌شود، مثلاً \lr{3.2} یعنی بخش دوم از فصل سوم.}

\section{زیربخش (\lr{subsection})}
\label{sec:chapters-subsection}

زیربخش، زیرمجموعه‌ی بخش است. برای ایجاد
زیربخش:

\begin{latin}
\begin{minted}{latex}
\subsection{|\pc{پیشینه}|}
\label{sec:background}
\end{minted}
\end{latin}

\subsection{سبک نمایش زیربخش}
\label{sec:chapters-subsection-style}

زیربخش‌ها با فونت \pnum{۱۳pt} و شماره‌ی آبی
کم‌رنگ نمایش داده می‌شوند.

\mnote{شماره‌ی زیربخش به صورت «فصل.بخش.زیربخش» نمایش داده می‌شود، مثلاً \lr{3.2.1}.}

\section{زیرزیربخش (\lr{subsubsection})}
\label{sec:chapters-subsubsection}

زیرزیربخش، عمیق‌ترین سطح عنوان‌بندی در
\lr{MatinBook} است. برای ایجاد زیرزیربخش:

\begin{latin}
\begin{minted}{latex}
\subsubsection{|\pc{مثال اول}|}
\label{sec:example-first}
\end{minted}
\end{latin}

\mnote{استفاده‌ی بیش از حد از زیرزیربخش، ساختار کتاب را پیچیده می‌کند. سعی کنید از سه سطح اول (فصل، بخش، زیربخش) استفاده کنید.}

\section{جدول مقایسه‌ی سطوح}
\label{sec:chapters-comparison}

جدول \cref{tab:heading-levels} مقایسه‌ی خلاصه‌ای
بین پنج سطح عنوان‌بندی را نشان می‌دهد.

\begin{matintable}{مقایسه‌ی سطوح عنوان‌بندی}{tab:heading-levels}
\begin{tblr}{
    width = \textwidth,
    colspec = {l X[c] X[c] X[c]},
    hlines, vlines,
    colsep = 6pt,
    rowsep = 4pt,
    row{1} = {font=\bfseries},
}
دستور & اندازه & شماره‌گذاری & در فهرست \\
\lr{\textbackslash part} & صفحه‌ی جداگانه & رومی & بله \\
\lr{\textbackslash chapter} & \pnum{۲۴/۲۰}pt & عربی & بله \\
\lr{\textbackslash section} & \pnum{۱۵}pt & فصل.بخش & بله \\
\lr{\textbackslash subsection} & \pnum{۱۳}pt & فصل.بخش.زیر & بله \\
\lr{\textbackslash subsubsection} & \pnum{۱۱}pt & فصل.بخش.زیر.زیر & بله \\
\end{tblr}
\end{matintable}

\section{ارجاع به فصل و بخش}
\label{sec:chapters-references}

برای ارجاع به فصل و بخش، از دستور \cmd{cref}
استفاده می‌کنیم. این دستور به‌طور خودکار نوع
مرجع را تشخیص می‌دهد و پیشوند مناسب (فصل،
بخش، زیربخش) را اضافه می‌کند:

\begin{latin}
\begin{minted}{latex}
% |\pc{ارجاع به فصل}|
\cref{chap:introduction}

% |\pc{ارجاع به بخش}|
\cref{sec:introduction}

% |\pc{ارجاع به چند مرجع}|
\cref{chap:introduction,chap:chapters}
\end{minted}
\end{latin}

\mnote{خروجی \lr{\textbackslash cref\{chap:introduction\}} می‌شود «فصل ۱» و خروجی \lr{\textbackslash cref\{sec:introduction\}} می‌شود «بخش ۱.۱».}

\subsection{نام‌گذاری برچسب‌ها}
\label{sec:chapters-label-naming}

برای اینکه برچسب‌ها معنادار باشند، از
پیشوندهای زیر استفاده کنید:

\begin{itemize}
    \item \textbf{فصل:} \lr{chap:xxx} —
    مثال: \lr{chap:introduction}.
    \item \textbf{بخش:} \lr{sec:xxx} —
    مثال: \lr{sec:motivation}.
    \item \textbf{زیربخش:} \lr{subsec:xxx} —
    مثال: \lr{subsec:history}.
    \item \textbf{زیرزیربخش:} \lr{subsubsec:xxx}.
\end{itemize}

\mnote{برچسب‌های معنادار، خواندن فایل‌های \lr{tex} را آسان‌تر می‌کنند و از خطاهای ارجاع جلوگیری می‌کنند.}

\section{مثال کامل: یک فصل با همه‌ی سطوح}
\label{sec:chapters-full-example}

در ادامه، یک فصل کامل با همه‌ی سطوح عنوان‌بندی
نمایش داده می‌شود:

\begin{latin}
\begin{minted}{latex}
\chapter{|\pc{مقدمه}|}
\label{chap:introduction}

\section{|\pc{انگیزه}|}
\label{sec:motivation}

|\pc{متن انگیزه...}|

\subsection{|\pc{پیشینه}|}
\label{subsec:background}

|\pc{متن پیشینه...}|

\subsubsection{|\pc{کارهای مرتبط}|}
\label{subsubsec:related}

|\pc{متن کارهای مرتبط...}|

\section{|\pc{تعریف مسئله}|}
\label{sec:problem}

|\pc{متن تعریف مسئله...}|

\section{|\pc{نتیجه‌گیری}|}
\label{sec:conclusion}

|\pc{متن نتیجه‌گیری...}|
\end{minted}
\end{latin}

\mnote{در این مثال، از پنج سطح عنوان‌بندی استفاده شده است. سعی کنید از چهار سطح اول استفاده کنید و از زیرزیربخش فقط در موارد ضروری.}

\section{خطاهای رایج}
\label{sec:chapters-mistakes}

\subsection{عنوان‌بندی بدون متن}
\label{sec:chapters-mistake-empty}

هرگز دو عنوان را پشت سر هم بدون متن قرار
ندهید. این کار، خواننده را گیج می‌کند:

\begin{latin}
\begin{minted}{latex}
% |\pc{اشتباه}|:
\section{|\pc{مقدمه}|}
\subsection{|\pc{پیشینه}|}

% |\pc{درست}|:
\section{|\pc{مقدمه}|}
|\pc{این فصل به بررسی... می‌پردازد.}|

\subsection{|\pc{پیشینه}|}
|\pc{پیشینه پژوهش...}|
\end{minted}
\end{latin}

\subsection{استفاده‌ی بیش از حد از زیرزیربخش}
\label{sec:chapters-mistake-deep}

اگر از زیرزیربخش بیش از حد استفاده کنید،
ساختار کتاب پیچیده می‌شود و خواننده مسیر
خود را گم می‌کند.

\textbf{راه‌حل:} از سه سطح اول (فصل، بخش،
زیربخش) استفاده کنید و فقط در موارد ضروری
به زیرزیربخش بروید.

\subsection{فراموش کردن \cmd{label}}
\label{sec:chapters-mistake-label}

اگر \cmd{label} را فراموش کنید، نمی‌توانید
به آن بخش ارجاع دهید و ممکن است در آینده
مشکل ایجاد شود.

\textbf{راه‌حل:} برای هر فصل و بخش مهم، یک
\cmd{label} معنادار تعریف کنید.

\section{تمرین‌ها}
\label{sec:chapters-exercises}

\begin{exercise}
    یک فصل با سه بخش و هر بخش با دو زیربخش
    بسازید. آیا ساختار در فهرست مطالب به‌درستی
    نمایش داده می‌شود؟
    \label{exc:chapters-structure}
\end{exercise}

\begin{exercise}
    به فصل و بخش‌های خود ارجاع دهید. آیا
    پیشوندها به‌درستی نمایش داده می‌شوند؟
    \label{exc:chapters-refs}
\end{exercise}

\begin{exercise}
    یک کتاب با دو بخش اصلی و چهار فصل بسازید.
    از دستور \cmd{part} برای ایجاد بخش‌های
    اصلی استفاده کنید.
    \label{exc:chapters-parts}
\end{exercise}

\begin{exercise}
    یک \cmd{chapter*} با \cmd{addcontentsline}
    بسازید. آیا در فهرست مطالب ظاهر می‌شود؟
    \label{exc:chapters-star}
\end{exercise}

\section{خلاصه}
\label{sec:chapters-summary}

در این فصل، با سیستم عنوان‌بندی \lr{MatinBook}
آشنا شدیم:

\begin{itemize}
    \item \lr{MatinBook} از پنج سطح عنوان‌بندی
    پشتیبانی می‌کند: \cmd{part}، \cmd{chapter}،
    \cmd{section}، \cmd{subsection}،
    \cmd{subsubsection}.

    \item مقیاس اندازه‌ها بر اساس استاندارد
    \lr{Boyer/Stewart} است: ۲۴/۲۰، ۱۵، ۱۳، ۱۱pt.

    \item برای فصل‌های بدون شماره، از
    \cmd{chapter*} با \cmd{addcontentsline}
    استفاده می‌شود.

    \item ارجاع به فصل و بخش با \cmd{cref}
    انجام می‌شود که پیشوند مناسب را به‌طور
    خودکار اضافه می‌کند.

    \item برچسب‌ها باید معنادار باشند و از
    پیشوندهای استاندارد (\lr{chap:}،
    \lr{sec:}، \lr{subsec:}) استفاده کنند.

    \item خطاهای رایج شامل عنوان‌بندی بدون
    متن، استفاده‌ی بیش از حد از زیرزیربخش،
    و فراموش کردن \cmd{label} است.
\end{itemize}

در فصل بعدی (\cref{chap:text})، با متن و
پاراگراف آشنا می‌شوید و یاد می‌گیرید که
چگونه متن فارسی و لاتین را با هم ترکیب کنید.